\ntsection{The Architect's Question}\label{the-architects-question}


As a RAG Q\&A workload scales from a single host to a fleet of serving
nodes, the architect must confront a new set of trade-offs: how does
aggregate throughput scale with added hardware, how does scheduling
latency affect user-perceived latency, and at what point does
autoscaling cease to be cost-effective? This chapter addresses these
questions through concrete arithmetic, not hand-waving.

\ntsection{Concept}\label{concept}


Fleet-level optimization treats the entire collection of serving hosts
as a single resource pool rather than independent islands. The central
insight is that aggregate system throughput is not merely the sum of
individual host throughputs; it is modulated by scheduling overhead,
network contention, and the cost of moving data between hosts and
storage. Two cost effects must be kept distinct. The \textbf{average}
cost per query declines as the fleet grows, because fixed costs (model
loading, infrastructure, staff) are amortized across more concurrent
requests. The \textbf{marginal} cost of each added host, however, is
roughly constant in this model --- one host buys \textasciitilde2.0
req/s at \textasciitilde\$20/hr, so the marginal cost stays near
\textasciitilde\$9.5 per req/s per hour (Ch 17 §7 frames the Nth host as
diminishing marginal \emph{benefit}, not falling marginal cost). The
economic lever is amortizing the fixed term, not the marginal one.

The key concepts are:

\begin{itemize}
\tightlist
\item
  \textbf{Aggregate throughput}: The total queries-per-second the fleet
  can sustain before latency SLA violations.
\item
  \textbf{Scheduling overhead}: The per-request cost of dispatching a
  query to an available host, including placement decisions and data
  movement.
\item
  \textbf{Autoscaling break-even}: The point at which adding another
  host reduces marginal cost per query versus the fixed cost of
  provisioning and maintaining that host.
\end{itemize}

\ntsection{Mental Model}\label{mental-model}


The mental model for fleet optimization replaces the single-host query
pipeline with a two-level hierarchy. At the lower level, each host
processes queries using its GPU/CPU pipeline. At the upper level, a
scheduler routes incoming requests across hosts, handling load
balancing, model parallelism where needed, and graceful degradation
under overload.

The scheduler's decisions have direct impact on the tail of the latency
distribution. Poor placement can cause some hosts to saturate while
others remain underutilized, inflating the 99th-percentile latency even
when aggregate throughput is ample. Conversely, intelligent
scheduling---such as consistent hashing for session affinity or
predictive placement based on historical patterns---keeps utilization
flat across the fleet and minimizes tail latency.

A critical assumption in this model is that the workload follows a known
arrival process. For the canonical scenario described here:
\textasciitilde2,000 registered users, \textasciitilde5\% actively using
the system concurrently (\textasciitilde100 concurrent \emph{users}, a
population measure), generating \textasciitilde10 rps average and
\textasciitilde40 rps peak. With \textasciitilde9,200 in-tokens and
\textasciitilde300 out-tokens per query, a \textasciitilde8.6 s service
time per request (1.08 s prefill + \textasciitilde7.5 s decode), and a
70B dense FP16 model running on 8×H100 (640 GB total HBM3), the fleet
must sustain \textbf{\textasciitilde40 requests/s at peak}, which --- by
Little's law --- corresponds to \textbf{λ·W ≈ 40 × 8.6 ≈ 344 requests in
flight}, not merely the \textasciitilde100 concurrent \emph{users} or
the nominal 40 rps. The distinction between ``concurrent users'' (a
population), the arrival rate (rps), and in-flight requests (λ·W) is
exactly what \hyperref[chap:17]{Chapter 17} develops; it is the difference between sizing
for a population and sizing a fleet against a latency bound.

\ntsection{Worked Example}\label{worked-example}


Take a fleet of N hosts, each with 8×H100 running a single 70B FP16
model instance. Two capacities must be reconciled ---
\textbf{throughput} (tokens/s) and \textbf{KV residency} (concurrent
requests) --- and neither can be reduced to a single ``QPS'' constant,
because QPS is workload-specific (per-request input/output tokens,
context length, batch).

\textbf{Latency floor first, then reconcile the two capacity bounds.} A
300-output-token request at a \textasciitilde25 ms/token decode rate
(\hyperref[chap:8]{Chapter 8}) takes \textasciitilde7.5 s of generation alone, plus 1.08 s
prefill, so the service time is \textasciitilde8.6 s per request (before
tool calls). Two separate ceilings bound how many requests a host can
\emph{serve per second}, and they must be reconciled --- they are not
interchangeable:

\begin{itemize}
\tightlist
\item
  \textbf{Throughput bound (tokens/s).} A host sustains an illustrative
  \textasciitilde35,000 tokens/s {[}ILLUSTRATIVE achievable range,
  \hyperref[chap:15]{Chapter 15}{]}. Dividing by the \textasciitilde9,500 tokens/request
  gives \textasciitilde3.7 requests/s --- \textbf{but this is not
  achievable here}, because it would require \textasciitilde3.7 × 8.6 ≈
  32 concurrent requests in flight, roughly double the host's
  KV-residency ceiling (C ≈ 18, \hyperref[chap:17]{Chapter 17}). The token-throughput figure
  is thus a non-binding upper edge for this canonical workload.
\item
  \textbf{KV-residency / latency bound (Little's law).} The real ceiling
  is set by concurrency and service time together: a host holds
  \textasciitilde17.5 concurrent full-context requests (KV ceiling,
  \hyperref[chap:17]{Chapter 17}) each for \textasciitilde8.6 s, so it serves C/W ≈ 17.5 /
  8.6 ≈ \textbf{2.0 requests/s} at steady state.
\end{itemize}

The binding per-host rate is the \emph{minimum} of the two --- here the
KV/latency bound. This is exactly why a token-throughput number must not
be divided by tokens-per-request to manufacture a ``QPS'': a token and a
request are different units, and prefill (compute-bound) versus decode
(bandwidth-bound) tokens are not fungible in that division. The per-host
request rate is a \emph{concurrency × service-time} quantity, not a
tokens ÷ tokens conversion.

\textbf{Call this quantity what it is.} The \textasciitilde2.0 req/s
figure is a \textbf{KV-residency / service-time analytical planning
bound under the canonical W assumption} --- not a measured serving
throughput. It is achievable only if serving C concurrent sequences does
\emph{not} inflate W beyond the assumed \textasciitilde8.6 s, and
neither a compute, bandwidth, nor scheduler constraint binds first. In
real continuous batching, W is a function of batch occupancy, scheduler
policy, prefill interference, decode batching, the arrival process,
tensor-parallel communication, and the context distribution --- so C and
W are not always independent constants. Treat \textasciitilde2.0 req/s
as an upper-bound model, and require benchmarked goodput (\hyperref[chap:14]{Chapter 14})
under the target arrival process for actual fleet sizing.

\textbf{Fleet aggregate throughput (illustrative form).} Keeping the
illustrative per-host QPS = Q for the algebra, aggregate throughput is

\[
\Lambda_\text{fleet} = N \times Q
\]

At the canonical peak of 40 rps, fleet size must satisfy \textbf{both}
the KV-residency and the throughput/latency bounds \emph{and whichever
binds first}. With the per-host rate at \textasciitilde2.0 req/s (the
binding Little bound above), a fleet for 40 rps starts at
\textasciitilde20 hosts on the Little-bound. The KV-residency bound
coincides here --- by Little's law, the in-flight count at 40 rps is λ·W
= 40 × 8.6 ≈ 344, and ⌈344/17.5⌉ ≈ 20 hosts --- so the two bounds are
\emph{co-binding} at the peak, exactly as \hyperref[chap:17]{Chapter 17} notes (the two
views are the same identity, λW/C = λ/(C/W)). The ``far fewer hosts
(\textasciitilde5)'' figure applies only to the 10-rps average, not the
40-rps peak. The correct statement is: the fleet must satisfy both, and
the \emph{larger} requirement wins.

\textbf{Scheduling overhead.} A lightweight round-robin scheduler adds
\textasciitilde0.5 ms of dispatch latency per request --- negligible
next to a real 7.5 s decode, but it must stay under \textasciitilde5\%
of the \emph{measured} per-request time.

\textbf{Autoscaling economics.} Each 8×H100 host costs
\textasciitilde\$20/hour (cloud, \$2.50/GPU-hr canonical) or
\textasciitilde\$8/hour amortized (owned) --- {[}ILLUSTRATIVE{]}. The
marginal cost per additional QPS is \$20 ÷ (QPS gained per host); the
economic lesson (diminishing cost of a marginal host once fixed costs
dominate) is sound, but the exact per-QPS figure is only meaningful
after the real per-host throughput is benchmarked.

\ntsection{Measurement}\label{measurement}


To validate fleet-level optimization claims, three core metrics should
be measured at regular intervals:

\begin{longtable}[]{@{}
  >{\raggedright\arraybackslash}p{0.2933\linewidth}
  >{\raggedright\arraybackslash}p{0.2933\linewidth}
  >{\raggedright\arraybackslash}p{0.2933\linewidth}@{}}
\toprule\noalign{}
\begin{minipage}[b]{\linewidth}\raggedright
Metric
\end{minipage} & \begin{minipage}[b]{\linewidth}\raggedright
Definition
\end{minipage} & \begin{minipage}[b]{\linewidth}\raggedright
Target
\end{minipage} \\
\midrule\noalign{}
\endhead
\bottomrule\noalign{}
\endlastfoot
\textbf{Aggregate QPS} & Total queries processed per second across all
hosts & ≥ 2× peak rps for comfort margin \\
\textbf{Scheduler-added p99 latency} & 99th-percentile time a request
waits in the fleet dispatcher (NOT end-to-end --- model execution is
multi-second here) & \textless{} 100 ms \\
\textbf{Scheduler-added median latency} & 50th-percentile dispatch
latency, scheduler overhead only & \textless{} 30 ms \\
\textbf{Scheduling overhead} & Fraction of per-query time spent in
dispatcher, not compute & \textless{} 5\% \\
\textbf{Utilization variance} & Standard deviation of per-host query
rate as a fraction of mean & \textless{} 0.2 \\

\label{tab:20.1}\end{longtable}

\emph{Note: end-to-end median/p99 latency is deliberately \textbf{not} a
fleet-level target for this canonical 9.5K-token workload --- a request
takes \textasciitilde8.6 s in service (1.08 s prefill +
\textasciitilde7.5 s decode), so an }end-to-end* \textless100 ms p99 is
infeasible (see the mini-case). The table therefore qualifies these as
scheduler-added (dispatch) latency targets only.*

Measurement methodology: use a distributed load generator (e.g., Locust
or a custom gRPC-driven pump) that sends queries at the target rps while
instrumenting each host's internal metrics---GPU occupancy, kernel
launch time, scheduler dispatch timestamp, and queue depth. Aggregate
these via a time-series database (InfluxDB or Prometheus) and compute
the above metrics over sliding 1-minute windows.

A practical pitfall: measuring only aggregate QPS without tracking
per-host variance can mask saturation in individual nodes. Always report
both the fleet-wide and per-host metrics.

\ntsection{Common Mistakes}\label{common-mistakes}


\begin{itemize}
\item
  \textbf{Treating hosts as independent}: Optimizing each host in
  isolation (e.g., tuning batch size per host) ignores the scheduling
  overhead that becomes dominant when the fleet is viewed as a unit. A
  suboptimal batch size on one host can increase dispatch frequency for
  the entire fleet.
\item
  \textbf{Overlooking data movement cost}: In RAG workloads, the
  \textasciitilde9,200 in-tokens and \textasciitilde300 out-tokens per
  query must be fetched from a vector database or cache. If every host
  pulls data independently, the network cost scales linearly with fleet
  size. A better approach is to colocate vector stores with serving
  hosts or use a shared cache layer that amortizes fetch cost across
  concurrent queries.
\item
  \textbf{Autoscaling too aggressively}: Adding hosts at the first sign
  of latency increase can be suboptimal if the bottleneck is scheduler
  throughput or cache contention, not compute capacity. Measure the
  marginal cost of adding a host before triggering a scale-out event.
\item
  \textbf{Ignoring cold-start cost}: When a host is freshly provisioned
  or re-warmed after idle, loading the full model weights into HBM takes
  30--120 seconds (a weight-load figure; distinct from the much shorter
  KV-cache warm-up of \textasciitilde2 s discussed in \hyperref[chap:15]{Chapter 15}, which
  re-populates only the retained KV state). Autoscaling policies must
  account for this weight-load warm-up window when scaling out a fresh
  host; otherwise the fleet experiences a spike in error rate or timeout
  during scale-out events.
\end{itemize}

\ntsection{Architecture Consequence}\label{architecture-consequence}


Fleet-level optimization reshapes the serving architecture in three
ways:

\begin{enumerate}
\def\labelenumi{\arabic{enumi}.}
\item
  \textbf{Centralized scheduler}: A lightweight service (e.g., based on
  etcd or a custom placement engine) maintains real-time view of
  per-host utilization and makes dispatch decisions. This service must
  be highly available and sub-millisecond in decision latency, otherwise
  it becomes the system bottleneck.
\item
  \textbf{Data colocation}: Vector indexes and model weights are stored
  close to the hosts that consume them. This may mean running a local
  Qdrant or Milvus instance per host, or using a shared object store
  (e.g., MinIO) with read-replicas per availability zone. The goal is to
  minimize the distance---measured in network hops---between the
  scheduler's placement decision and the data fetch.
\item
  \textbf{Graceful degradation}: The fleet should have a mode where,
  under extreme overload, it serves stale or truncated answers rather
  than failing entirely. This could mean returning the top-k results
  from a cached embedding rather than recomputing, or surfacing a ``try
  again'' message with an estimated wait time. The architectural
  contract must explicitize this trade-off so product teams can set
  appropriate SLOs.
\end{enumerate}

\ntsection{What We Still Don't Know}\label{what-we-still-dont-know}


Several open questions remain at the fleet level:

\begin{itemize}
\item
  \textbf{Heterogeneous hardware}: What happens when hosts mix H100,
  A100, and CPU-only nodes? The scheduler must handle different
  throughput ceilings and memory capacities, and the aggregate
  throughput calculation becomes a weighted sum rather than a simple
  multiple. Real-world measurements of mixed-device fleets are scarce;
  most published benchmarking uses homogeneous hardware, so the
  crossover point where heterogeneity hurts versus helps is not well
  quantified.
\item
  \textbf{Multi-model serving}: When multiple model sizes coexist (e.g.,
  70B for complex queries, 7B for quick ones), the scheduling problem
  gains a dimensionality dimension---placement must balance model size
  against query urgency and host memory. Mixed-precision serving (FP16,
  FP8, INT4) adds another layer: lower precision increases throughput
  but may degrade answer quality on subtle reasoning tasks, and the
  quality-throughput trade-off curve is model-specific.
\item
  \textbf{Economic auto-scaling}: The break-even point between owning
  hardware versus cloud-on-demand depends on utilization patterns that
  are workload-specific. A fleet that runs at 20\% utilization on
  average may be cheaper in the cloud; one at 80\% utilization favors
  owned hardware. The transition point is a function of capital
  expenditure amortization period, electricity cost, and the volatility
  of user demand. Published TCO analyses typically assume steady
  utilization, but bursty workloads with short-lived spikes can make
  ownership risky unless over-provisioned.
\item
  \textbf{Tail-latency arbitration}: When two queries compete for the
  same GPU HBM bandwidth, the 99th-percentile latency can spike
  disproportionately. We lack a general-purpose scheduler policy that
  guarantees a target tail latency under arbitrary arrival patterns;
  this remains an area of active research. Some systems use priority
  queues or token buckets to cap per-query bandwidth, but these
  mechanisms add scheduling complexity and can increase median latency.
\item
  \textbf{Cross-query caching effects}: In RAG workloads, the
  \textasciitilde9,200 in-tokens and \textasciitilde300 out-tokens per
  query overlap significantly across users asking related questions. The
  potential for caching intermediate embeddings or partial KV caches is
  substantial, but the cost of cache management (invalidation,
  coherence, metadata overhead) must be quantified against the compute
  savings. We lack large-scale measurements of caching efficiency across
  diverse query populations.
\end{itemize}

\section{End-of-Chapter Mini-Case: Aggregate QPS, Marginal Cost, and an
Infeasible
SLO}\label{end-of-chapter-mini-case-aggregate-qps-marginal-cost-and-an-infeasible-slo}

\textbf{Scenario}: A RAG Q\&A fleet serves 2,000 users with the
canonical parameters: 5\% concurrent, 10 rps average / 40 rps peak,
\textasciitilde9,200 in-tokens + \textasciitilde300 out-tokens per
query, 70B dense FP16 model, 8×H100 per host, 8 hosts total.

\textbf{Question}: What is the fleet's aggregate QPS capacity under the
canonical KV/latency bound, and what is the marginal cost per QPS when
scaling from 8 to 16 hosts? Is an end-to-end 99th-percentile latency
under 100 ms even feasible for this workload --- and if not, what can
the scheduler bound?

\textbf{Solution} (illustrative algebra --- the per-host figure below is
an {[}ILLUSTRATIVE{]} estimate for showing the fleet math; the real
per-host QPS must come from a deployment benchmark, \hyperref[chap:14]{Chapter 14}, and is
bounded by KV residency per \hyperref[chap:17]{Chapter 17}):

\begin{itemize}
\tightlist
\item
  \textbf{Per-host request rate --- the KV/latency bound (Ch. 17
  basis).} A host holds \textasciitilde17.5 concurrent full-context
  requests (the KV-residency ceiling, \hyperref[chap:17]{Chapter 17}; byte-accurate 17.51,
  conservative integer ceiling 17) and each request is in service
  \textasciitilde8.6 s (1.08 s prefill + \textasciitilde7.5 s decode of
  300 output tokens at \textasciitilde25 ms/token, \hyperref[chap:8]{Chapter 8}). So a host
  serves C/W ≈ 17.5 / 8.6 ≈ \textbf{2.0 requests/s} at steady state.
  This is the defensible capacity number: it is derived from concurrency
  × service time (Little's law), not from dividing a token throughput by
  tokens per request. (Do NOT divide \textasciitilde35,000 tokens/s by
  \textasciitilde9,500 tokens/request to get \textasciitilde3.7 req/s
  --- that would require \textasciitilde32 concurrent requests, which
  the \textasciitilde17.5-request KV ceiling cannot hold.)
\item
  \textbf{Fleet aggregate (8 hosts):} 8 × 2.0 ≈ \textbf{16 req/s}, which
  is \textbf{below} the canonical 40-rps peak. On this binding bound the
  idealized minimum is ⌈40 / 2.0⌉ = \textbf{20 hosts}; allowing
  \textasciitilde10\% scheduling loss raises this floor to 40 / (2.0 ×
  0.9) ≈ 22.2 → \textbf{23 hosts}. (The KV-residency bound, \hyperref[chap:17]{Chapter 17},
  is satisfied at far fewer hosts --- \textasciitilde5 at the 10-rps
  average --- but the throughput/latency bound binds here, so the larger
  requirement governs.)
\item
  \textbf{Marginal cost (real basis):} adding one host at
  \textasciitilde\$20/hr (cloud, {[}ILLUSTRATIVE{]}) buys
  \textasciitilde2.0 req/s, so marginal cost ≈ \$20 ÷ 2.0 = \textbf{\$10
  per req/s per hour.} (The earlier illustrative \$0.012 figure assumed
  250 req/s per host --- a \textasciitilde120× understatement that comes
  entirely from the impossible latency; with a real multi-second service
  time per request the per-host rate is \textasciitilde2.0, not 250.)
\item
  \textbf{Algebra practice (kept {[}ILLUSTRATIVE{]}):} the per-host 250
  QPS / fleet 2,000 QPS / \$0.012 per QPS arithmetic from the original
  exercise is retained only to show the \emph{structure} of the fleet
  math (marginal host cost ÷ marginal capacity). It deliberately uses a
  hypothetical per-host QPS that is not achievable for this
  9.2K-in/300-out workload; the architect must substitute the measured
  per-host rate from a \hyperref[chap:14]{Chapter 14} deployment benchmark, or the
  \textasciitilde2.0 req/s bound derived here.
\item
  \textbf{The 100-ms end-to-end p99 is infeasible under the stated
  workload, and that is the honest answer.} A request here takes
  \textasciitilde7.5 s of decode alone (300 output tokens ×
  \textasciitilde25 ms/token, \hyperref[chap:8]{Chapter 8}), so no scheduler-overhead
  percentage can bring end-to-end p99 under 100 ms --- the model
  execution already exceeds the budget by two orders of magnitude. What
  the scheduler \emph{can} bound is the overhead it adds \emph{on top
  of} model execution: a lightweight round-robin dispatcher adds
  \textasciitilde0.5 ms per request, which is negligible against a
  multi-second decode, and a \textless=5\% scheduler-added p99 is
  comfortably achievable. The architect's job here is to reject the
  impossible SLO and state what is actually controllable, rather than
  tune arithmetic until a 100-ms figure appears reachable.
\end{itemize}

\begin{figure}
\centering
\pandocbounded{\includegraphics[keepaspectratio,alt={Fleet capacity (KV/service-time analytical bound, \textasciitilde2.0 req/s/host) and per-host offered-load utilization, vs.~host count at the canonical peak load (40 rps). Top panel: capacity against host count from the KV/service-time analytical bound (\textasciitilde2.0 req/s/host), with a \textasciitilde95\% scheduling-overhead line. Bottom panel: per-host offered-load utilization (ρ = λ/(hosts·bound)) at the canonical 40 rps peak --- this is an analytical-capacity utilization, NOT GPU utilization. This is the honest quantity for the tail: p99 is driven by per-host utilization, not by host count alone. Adding hosts lowers per-host utilization (more headroom) hence lowers the expected tail; the 70\%-utilization and 100\%-saturation reference lines show the provisioning target. The tail also depends on the arrival process, scheduling policy, batching, and the service-time distribution, which a provisioning model alone does not capture. The \textasciitilde2.0 req/s line is the KV/service-time analytical bound, not a benchmarked serving throughput. {[}ILLUSTRATIVE{]}{[}DERIVED{]}},width=\maxwidth{\textwidth},keepaspectratio]{figures/fig-20-2001.pdf}}
\caption[Fleet capacity (KV/service-time analytical bound, \textasciitilde2.0 req/s/host)…]{Fleet capacity (KV/service-time analytical bound, \textasciitilde2.0 req/s/host) and per-host offered-load utilization, vs.~host count at the canonical peak load (40 rps). Top panel: capacity against host count from the KV/service-time analytical bound (\textasciitilde2.0 req/s/host), with a \textasciitilde95\% scheduling-overhead line. Bottom panel: per-host \emph{offered-load} utilization (ρ = λ/(hosts·bound)) at the canonical 40 rps peak --- this is an \emph{analytical-capacity} utilization, NOT GPU utilization. This is the honest quantity for the tail: p99 is driven by per-host utilization, not by host count alone. Adding hosts lowers per-host utilization (more headroom) hence lowers the expected tail; the 70\%-utilization and 100\%-saturation reference lines show the provisioning target. The tail also depends on the arrival process, scheduling policy, batching, and the service-time distribution, which a provisioning model alone does not capture. The \textasciitilde2.0 req/s line is the KV/service-time analytical bound, not a benchmarked serving throughput. {[}ILLUSTRATIVE{]}{[}DERIVED{]}}

\label{fig:20.1}\end{figure}
